\emph{(i).} $\alpha$ is a $*$-automorphism with $\alpha^2 = \mathrm{id}$, so $E = \tfrac12(\mathrm{id} + \alpha)$ is a positive unital projection onto $\Aev$ fixing $\Aev$ pointwise and killing $\Aod$ (Definition \claimref{obj.singularity.grading-parity-automorphism}). For a state $\omega$, $\omega \circ E = \tfrac12(\omega + \omega \circ \alpha)$ is a convex combination of states, hence a state; it is $\alpha$-invariant because $E \circ \alpha = E$, and it equals $\omega$ on $\Aev$ because $E$ fixes $\Aev$. If $\omega'$ is even and agrees with $\omega$ on $\Aev$, then for every $x$, $\omega'(x) = \omega'(Ex) = \omega(Ex)$, since $Ex \in \Aev$ and an even state satisfies $\omega' = \omega' \circ E$; so $\omega' = \omega \circ E$. Restriction to $\Aev$ and $\eta \mapsto \eta \circ E$ are therefore inverse bijections between the states of $\Aev$ and the even states of $\Aa$. The observable clause of Fact \claimref{singularity.imported-univalence-superselection-rule} makes $\Aev$ the observables; every even element $x$ has $\omega(x) = (\omega \circ E)(x)$, so no observable separates a state from its even part, and (P), $\omega \circ \alpha = \omega$, is the statement $\omega = \omega \circ E$. On one mode $v_k \in \Aod$, so $\omega(v_k)$, the off-diagonal entry of the ascended marginal in the number basis, is among the values (P) fixes. The check \texttt{check\_even\_part\_is\_unique\_even\_extension} verifies on two modes, for a non-even positive $\varrho$, that $\tfrac12(\varrho + \Gamma\varrho\Gamma)$ is positive of trace one, agrees with $\varrho$ on all eight even matrix units, vanishes on the eight odd ones, and is $\Gamma$-invariant, while $\varrho$ itself has a nonzero odd value.

\emph{The flow at imaginary time.} On one mode take the number basis, $v = E_{01}$, $v^{*} = E_{10}$, $v^{*}v = N_k = E_{11}$, and the Killing flow $\alpha_t = \operatorname{Ad}(e^{it\omega_k N_k})$ (the generator differs from $\omega_k K_k$ by a multiple of the identity, which does not change $\operatorname{Ad}$). Then $\alpha_{ip}(E_{ij}) = e^{-p\omega_k N_k} E_{ij} e^{p\omega_k N_k} = r^{\,j - i} E_{ij}$ with $r := e^{p\omega_k} > 0$, and for matrices the analytic continuation in (R) is this substitution, so the KMS relation for a pair $(a, b)$ reads $\omega(a\,\alpha_{ip}(b)) = \omega(b\,a)$.

\emph{(ii).} For $\varrho = \begin{psmallmatrix} 1-n & c \\ c & n\end{psmallmatrix}$, the four ladder pairs give: $(v^{*}, v)$: $r\,\omega(v^{*}v) = \omega(vv^{*})$, that is $rn = 1 - n$, which is $n = (r+1)^{-1}$; $(v, v^{*})$: $r^{-1}\omega(vv^{*}) = \omega(v^{*}v)$, the same equation; $(v, v)$ and $(v^{*}, v^{*})$: both sides vanish since $v^2 = 0$. The value $c$ enters none of them. $\varrho$ is positive exactly when $c^2 \leq n(1-n)$; its pinching is $\operatorname{diag}(1-n, n)$, the Gibbs state, as Theorem \claimref{singularity.descended-form-state-independent} says it must be for any state; and $\omega(v) = \operatorname{tr}(\varrho E_{01}) = c$. The check \texttt{check\_ladder\_kms\_admits\_coherence} verifies all of this for five values of $r$ (including $r = 1$) and the admissible $c$ among three, and that the same states fail the relation on the whole factor.

\emph{(iii).} Impose the relation for all sixteen pairs of matrix units on an unknown $\varrho = \begin{psmallmatrix} a & b \\ c & d \end{psmallmatrix}$. The pair $(E_{00}, E_{01})$ gives $r\,\omega(E_{01}) = \omega(E_{01}E_{00}) = 0$, so $\omega(E_{01}) = c = 0$ (the entry $\operatorname{tr}(\varrho E_{01})$ is $c$); likewise $(E_{11}, E_{10})$ gives $b = 0$; the pair $(E_{10}, E_{01})$ gives $r\,d = a$. With $a + d = 1$ the solution is unique: $\varrho = \operatorname{diag}(r, 1)/(r + 1)$, the Gibbs state with occupation $(r+1)^{-1}$, which is diagonal, so $\omega(v_k) = 0$. This is the one-mode case of the statement that Gibbs states of even generators are even (filed alongside, synthesis.univalence-from-even-stationary-dynamics, clause (2)); it is computed here because the zero-weight case is needed in the form of traciality. At zero weight $r = 1$ the relation is $\omega(ab) = \omega(ba)$ for all $a, b$, whose only normalised solution on $\Mat_2(\C)$ is the trace state $\tfrac12\operatorname{tr}$, again diagonal, with occupation $\tfrac12$. The check \texttt{check\_full\_kms\_forces\_single\_mode\_evenness} assembles the sixteen linear relations plus normalisation, verifies that they have rank four (a unique solution) and that the Gibbs state solves them, for the same five values of $r$.

\emph{(iv).} In the Jordan-Wigner representation on $\C^2 \otimes \C^2$, $v_1 = v \otimes 1$, $v_2 = P \otimes v$ with $P = \operatorname{diag}(1, -1) = 1 - 2N$, $\Gamma = P \otimes P$; then $\{v_1, v_2\} = \{v_1, v_2^{*}\} = 0$ and $\{v_1, v_1^{*}\} = 1$, and $1 - 2v_2^{*}v_2 = 1 \otimes P$. The element $y := (v_1 + v_1^{*})(1 - 2v_2^{*}v_2) = X \otimes P$, $X = v + v^{*}$, is odd: $\Gamma y \Gamma = (PXP) \otimes P = -y$. Put $\varrho = \varrho_1 \otimes \varrho_2 + c\,y$. Since $\operatorname{tr}(X) = \operatorname{tr}(P) = \operatorname{tr}(XP) = 0$, the term $c\,y$ contributes nothing to the expectation of any $E_{ij} \otimes 1$ (the mode-1 factor) or of $1$, $1 \otimes N$, $v_2 = P \otimes v$, $v_2^{*}$ (the mode-2 factor), so the two marginals are $\varrho_1$ and $\varrho_2$; and $\operatorname{tr}(\varrho\, y) = c\operatorname{tr}(y^2) = 4c$. In the basis ordered by mode 2, $\varrho$ is a direct sum of two $2 \times 2$ blocks $\begin{psmallmatrix} (1-n_1)g & \pm c \\ \pm c & n_1 g\end{psmallmatrix}$, $g$ an eigenvalue of $\varrho_2$, so $\varrho \geq 0$ exactly when $c^2 \leq n_1(1 - n_1)\min(n_2, 1 - n_2)^2$, which leaves room for $c \neq 0$. The check \texttt{check\_cross\_mode\_odd\_correlation\_with\_even\_gibbs\_marginals} verifies the CAR relations, the oddness of $y$, positivity, the marginals and the value $4c$ for three pairs of Gibbs marginals.

\emph{(v).} Let $x$ be odd in the algebra of one set of modes and $y$ odd in that of a disjoint set. Odd elements of disjoint modes anticommute ($v_j$ and $v_k^{(*)}$ do for $j \neq k$, and products of an odd number of them inherit it), so $xy = -yx$. A product state would give $\omega(xy) = \omega(x)\omega(y)$ and $\omega(yx) = \omega(y)\omega(x)$, and $\omega(xy) = -\omega(yx)$ then forces $\omega(x)\omega(y) = 0$. If both restrictions were non-even there would be odd $x$, $y$ with $\omega(x)\omega(y) \neq 0$, a contradiction. One non-even factor is possible: with $\varrho_1$ as in (ii) and $\varrho_2$ even, $\varrho_1 \otimes \varrho_2$ is a state, product across the two modes. The check \texttt{check\_graded\_product\_has\_one\_odd\_factor\_at\_most} verifies $xy = -yx$ for $x = X \otimes 1$, $y = P \otimes X$; that the one-odd-factor state $\varrho_1 \otimes \varrho_2$ (with $\varrho_2$ even) satisfies $\omega(ab) = \omega(a)\omega(b)$ for all sixteen pairs with $a \in \{E_{ij} \otimes 1\}$ and $b \in \{1 \otimes E_{kk},\, P \otimes E_{kl}\}$; and that for two coherent factors the Jordan-Wigner tensor product $\varrho_1 \otimes \varrho_2'$ is a state that is not a product: $\omega(xy) = 0$ while $\omega(x)\omega(y) = 2c_1(1 - 2n_1)\,2c_2 \neq 0$ for $n_1 \neq \tfrac12$.

\emph{Conclusion.} (P) is $\omega(z) = 0$ for every odd $z$. The odd elements of a single mode factor are spanned by $v_k, v_k^{*}$, so (P) contains the single-mode part, supplied by (iii) under the full-factor reading of (R) and not by (R) as stated, by (ii); the remaining odd elements involve two or more modes, and (iv) shows their expectations are not fixed by any condition on single-mode marginals, while (i) shows none of this is visible to an even observable.
